% Compile: pdflatex (two passes for cross-refs).  Engine: MiKTeX pdflatex.
\documentclass[11pt]{article}

\usepackage[letterpaper,margin=1in]{geometry}
\usepackage{mathptmx}            % Times text + math (NBER house look), base-safe
\usepackage{amsmath,amssymb}
\usepackage{booktabs}
\usepackage{array}
\usepackage{graphicx}            % \resizebox for one wide display
\usepackage{caption}
\usepackage{microtype}
\usepackage[hidelinks]{hyperref}

\captionsetup{font=small,labelfont=bf,skip=4pt}
\setlength{\parskip}{2pt}
\setlength{\emergencystretch}{2em}
\renewcommand{\arraystretch}{1.05}

% compact section heading spacing
\makeatletter
\renewcommand\section{\@startsection{section}{1}{\z@}%
  {-2.2ex \@plus -1ex \@minus -.2ex}{1.0ex \@plus .2ex}%
  {\normalfont\large\bfseries}}
\renewcommand\subsection{\@startsection{subsection}{2}{\z@}%
  {-1.8ex \@plus -1ex \@minus -.2ex}{0.6ex \@plus .2ex}%
  {\normalfont\normalsize\bfseries}}
\makeatother

\newcommand{\Sct}{S_{ct}}

\title{\vspace{-1.2cm}\bfseries An Empirical Framework for the Health-Cost
Incidence, Persistence, and Distributional Burden of Climate Shocks}
\author{Ahwaz Akhtar\\[-2pt]\normalsize GWU}
\date{\normalsize June 2026}

\begin{document}
\maketitle
\vspace{-0.7cm}

%================================================================
\section{Setting and identification}
%================================================================
The analysis uses two panels that share a common county-level climate-shock
series but measure different outcome domains. \emph{Panel~A} is a
county~$\times$~year panel (demand side): $\sim$41{,}400 county-years, 3{,}155
counties, 2011--2023, with outcomes for individual-market benchmark premiums,
medical-debt share, real personal income per capita, and civilian employment.
\emph{Panel~B} is a hospital~$\times$~year panel (supply side): $\sim$59{,}900
hospital-years, 5{,}119 hospitals (CMS Certification Numbers), with operating and
net margins and uncompensated care. Each hospital is geo-coded to a single county
(modal county per ZIP, facility location) and inherits that county's shocks, so
$S_{it}\equiv S_{c(i),t}$. Both panels carry county (or hospital) and year fixed
effects and cluster standard errors at the state level ($\sim$50 clusters).

All papers use four binary shocks defined at the county-year level:
$\textsf{Is\_Extreme\_Drought}$ (Palmer PDSI $\le-4$);
$\textsf{High\_CDD}$ / $\textsf{High\_HDD}$ (cooling/heating degree-days in the top
quartile of the county's own historical distribution); and
$\textsf{High\_AQI\_Max}$ (annual maximum AQI $>100$). Generically
$\Sct\in\{0,1\}$. These are \emph{recurring} shocks, not one-time absorbing
treatments: a county may enter, leave, and re-enter drought. This rules out
canonical staggered-adoption DiD as the \emph{primary} design and motivates the
persistence estimators of Section~3; Section~2 adds a first-onset cohort DiD as a
contamination-free cross-check. Defining shocks against each county's own
1990--2000 baseline distribution makes ``shock'' an unusual draw relative to local
historical norms rather than an absolute threshold.

%================================================================
\section{Paper 1 --- Incidence: do climate shocks raise health costs?}
%================================================================
\subsection{Distributed-lag two-way fixed effects}
Paper~1 asks whether contemporaneous and lagged exposure shifts the \emph{level} of
the outcome. For outcome $Y$ and shock $S$,
\begin{equation}
Y_{ct}=\alpha_c+\gamma_t+\sum_{k=0}^{2}\beta_k\,S_{c,t-k}
        +\mathbf{X}_{ct}'\delta+\varepsilon_{ct},
\label{eq:dl}
\end{equation}
with county FE $\alpha_c$, year FE $\gamma_t$, and time-varying controls
$\mathbf{X}_{ct}$ (log real income per capita, unemployment, Medicaid-expansion
status); the hospital version replaces $c$ with hospital $i$. Equation~\eqref{eq:dl}
is estimated \emph{once per (outcome, shock) pair} on the full panel; the three lag
terms enter the same regression. The subscripts run $t,t-1,t-2$ because a shock can
only affect current and \emph{subsequent} outcomes: a single drought in year $t$
appears as the $\beta_0$ regressor for the year-$t$ outcome, the $\beta_1$ regressor
for the year-$(t{+}1)$ outcome, and the $\beta_2$ regressor for year $t{+}2$. The
identifying variation is counties that \emph{switch} shock status; always- and
never-shocked counties are absorbed by $\alpha_c$. The cumulative (impulse-response)
effect of one shock-year traced across all three horizons is
\begin{equation}
\hat\beta_{\mathrm{cum}}=\hat\beta_0+\hat\beta_1+\hat\beta_2,
\qquad
\widehat{\operatorname{Var}}(\hat\beta_{\mathrm{cum}})
   =\sum_{j=0}^{2}\sum_{k=0}^{2}\widehat{\operatorname{Cov}}(\hat\beta_j,\hat\beta_k),
\label{eq:cum}
\end{equation}
the off-diagonal covariances taken from the state-clustered vcov via a Wald
linear-combination test. On the hospital panel, heat raises uncompensated care
($\textsf{High\_CDD}$: $+1.22$ pp of net patient revenue, $p=0.063$; cold:
$+\$1.55$M per hospital, $p=0.022$). The negative drought coefficient on
uncompensated care ($-1.68$ pp, $p=0.032$) is a measurement artifact---billing and
credit-bureau measures undercount distress in drought-stricken, lower-utilization
areas---and is flagged throughout rather than read causally.

\subsection{A sharper control group: the 2012 drought natural experiment}
Because the lag model's control group for any drought year silently includes
counties droughted in \emph{other} years, we add a first-onset cohort DiD. Let
$\text{Treated}_c=1$ if county $c$'s first-ever extreme-drought year in the panel is
2012, and $\text{Post}_t=\mathbf 1[t\ge2012]$:
\begin{equation}
Y_{ct}=\alpha_c+\gamma_t+\tau\,(\text{Treated}_c\times\text{Post}_t)+\varepsilon_{ct}.
\label{eq:did}
\end{equation}
The FEs absorb both main effects, so $\hat\tau$ is the ATT: 139 first-onset-2012
counties against 2{,}534 \emph{never}-droughted controls (counties first droughted in
other years are excluded, which is what cleans the control group). With a single
pre-period (2011) the regression returns the $2\times2$ difference exactly.

The sharp $2\times2$ assumes \emph{unconditional} parallel trends, which is strong
here: treated counties skew rural-agricultural (Georgia, the Mountain West, the
Plains) while controls skew urban and coastal. We therefore also report the
Sant'Anna--Zhao (2020) doubly-robust ATT, consistent if \emph{either} the propensity
or the outcome-regression model is correct, conditioning only on \emph{baseline}
(2011) covariates---log population, baseline median household income, Census
division---and deliberately not on contemporaneous income, employment, or premiums,
which are post-treatment mediators (bad controls):
\begin{equation}
\resizebox{0.96\hsize}{!}{$\displaystyle
\widehat\tau_{\mathrm{dr}}=\mathbb E\!\left[\left(\frac{D}{\mathbb E[D]}
-\frac{\hat\pi(X)\,(1-D)/(1-\hat\pi(X))}
       {\mathbb E\!\left[\hat\pi(X)(1-D)/(1-\hat\pi(X))\right]}\right)
\bigl(\Delta Y-\widehat m_0(X)\bigr)\right]$}
\label{eq:drdid}
\end{equation}
where $\Delta Y$ is the pre-to-post outcome change (post collapsed to the 2012--2023
mean), $\hat\pi(X)$ the propensity score, and $\widehat m_0(X)$ the control
outcome-change regression. Results are in Table~\ref{tab:incidence}.

\begin{table}[h]
\centering
\caption{Incidence: 2012-drought natural experiment (county panel, \$2023)}
\label{tab:incidence}
\small
\begin{tabular}{lrrc}
\toprule
Outcome & Uncond.\ $2\times2$ $\hat\tau$ & Doubly-robust $\hat\tau_{\mathrm{dr}}$ & DR 95\% CI \\
\midrule
Personal income per capita & $-1{,}311$ & $\mathbf{-1{,}451}$ & $[-2{,}461,\,-441]$ \\
Civilian employed (persons) & $-2{,}053$ & $-871$ & $[-1{,}719,\,-23]$ \\
Median household income & $-992$ & $-1{,}186$ & $[-2{,}140,\,-232]$ \\
Medical-debt share & $-0.006$\,(n.s.) & $-0.011$ & $[-0.019,\,-0.002]$ \\
\bottomrule
\end{tabular}

\vspace{2pt}
\footnotesize\raggedright
139 treated (first-onset 2012) vs.\ 2{,}534 never-exposed; state-clustered.
$2\times2$: income $z=-2.28$ ($p=0.027$), employment $z=-4.29$ ($p=0.0001$).
Income is robust to conditioning (if anything stronger); employment attenuates
$\sim$58\% and is the composition-sensitive secondary finding. Medical debt is
measurement-fragile. Re-estimated as a full Callaway--Sant'Anna model across
\emph{all} drought cohorts the average is null (PCPI $+\$350$, s.e.\ 585), so the
negative real-economy effect is specific to the large, clean 2012 event and is an
intent-to-treat ``effect of first onset.''
\end{table}

The unconditional $2\times2$ also leans on \emph{few treated clusters}: 67\% of the 139
treated counties sit in four states (GA, CO, NE, NM; 17 states in all), so state-clustered
analytic $p$-values risk over-rejection when treatment variation lives in few clusters. A
wild cluster bootstrap-$t$ (Cameron--Gelbach--Miller 2008, Webb weights, null imposed,
$B=9{,}999$) and Fisher randomization inference (2{,}000 re-draws of the treated labels),
both computed on the Frisch--Waugh--Lovell-residualized model so the bootstrap need not
carry 3{,}155 absorbed effects, confirm the income headline: $p_{\mathrm{wcb}}=0.036$ and
$p_{\mathrm{RI}}=0.007$ (analytic $0.027$), with a bootstrap CI $[-2{,}911,\,-138]$ that
excludes zero. Employment likewise clears the few-cluster bar ($p_{\mathrm{wcb}}=0.003$),
so its fragility is one of \emph{conditioning} (above), not cluster count; the
medical-debt-share null is unchanged. Implementation:
\texttt{did\_robustness/01\_wild\_cluster\_bootstrap.R}.

Because the county panel begins in 2011, the 2012 cohort has a single pre-period and
parallel trends is untestable within the panel. BEA CAINC1 per-capita income, however,
runs back to 1990, giving \emph{21} pre-treatment years for the same 139 treated and
2{,}534 never-exposed counties. Over 1990--2011 the treated-vs-control differential
\emph{linear} trend---the drift that would bias a DiD by extrapolating into the
post-period---is $-\$69$ per year (s.e.\ 89, $p=0.44$), i.e.\ economically and
statistically flat. A year-by-year event study (relative to 2011) does reject the joint
null of no pre-treatment gaps ($F=6.9$, $p<0.001$), but this reflects modest, business-cycle-
correlated wiggles between rural-agricultural (treated) and urban (control) counties rather
than a secular divergence---precisely the compositional difference the doubly-robust
estimator conditions on, and which it shows only \emph{strengthens} the income effect
($-\$1{,}451$). The two-decade trajectories move together before 2012 and visibly diverge
after (Fig.\ in \texttt{Analysis/did/robustness/bea\_pretrends\_1990\_2011.png}); the flat
linear pre-trend is the corroborating defense the single-pre-period design previously lacked.
Implementation: \texttt{did\_robustness/05\_bea\_pretrends\_1990\_2011.R}.

%================================================================
\section{Paper 2 --- Persistence: scarring, reversibility, dose}
%================================================================
\subsection{Onset/persist/exit local projections}
Converting the binary series into a transition state from $(\Sct,S_{c,t-1})$ gives
\emph{onset} $D^{on}\!=\!(1,0)$, \emph{persist} $D^{per}\!=\!(1,1)$, and \emph{exit}
$D^{ex}\!=\!(0,1)$, with $(0,0)$ omitted as the reference. Following Jord\`a (2005),
a separate regression is run per horizon $h$:
\begin{equation}
Y_{i,t+h}=\alpha_i+\gamma_t
   +\beta^{on}_h D^{on}_{it}+\beta^{per}_h D^{per}_{it}+\beta^{ex}_h D^{ex}_{it}
   +\varepsilon_{i,t+h},\qquad h=0,1,2,
\label{eq:lp}
\end{equation}
where the LHS is the outcome $h$ years after the transition year $t$ and the dummies
are always indexed by $t$. Each coefficient is measured against the never-shocked
$(0,0)$ reference: $\beta^{on}_h$ is the entry effect, $\beta^{per}_h$ the interior
of a sustained spell, and $\beta^{ex}_h$ the first post-shock year. If a shock were
fully reversible the entry cost would be exactly undone on exit, so the null and the
asymmetry statistic are
\begin{equation}
H_0:\ \beta^{on}_h+\beta^{ex}_h=0,
\qquad
\hat L_h=\hat\beta^{on}_h+\hat\beta^{ex}_h,
\qquad
z_h=\frac{\hat L_h}{\sqrt{V_{11}+V_{33}+2V_{13}}}\xrightarrow{d}\mathcal N(0,1),
\label{eq:sym}
\end{equation}
the covariance $V_{13}$ taken directly from the single fitted model's vcov.
A nonzero $\hat L_h$ means exit does not return the unit to its pre-shock
trajectory---scarring.

\subsection{Cumulative dose}
A second lens asks whether the \emph{stock} of past exposure matters. Define a
never-resetting counter $C_{it}=\sum_{s\le t}S_{is}$ and estimate
\begin{equation}
Y_{it}=\alpha_i+\gamma_t+\beta_1 C_{it}+\beta_2 C_{it}^2+\varepsilon_{it},
\qquad
\frac{\partial Y}{\partial C}\Big|_{C=x}=\beta_1+2\beta_2 x,
\label{eq:dose}
\end{equation}
with a binned variant $\sum_b\delta_b\mathbf 1[C_{it}\in b]$ over
$b\in\{1\text{--}3,4\text{--}6,7\text{--}9,10{+}\}$ (reference $C=0$) and a
high-vs-low contrast $\hat\Delta=\hat\delta_{10+}-\hat\delta_{1\text{--}3}$.
Table~\ref{tab:persist} reports drought $\to$ hospital operating margin.

\begin{table}[h]
\centering
\caption{Persistence: extreme drought $\to$ hospital operating margin ($N=57{,}302$)}
\label{tab:persist}
\small
\begin{tabular}{lrl}
\toprule
Quantity & Estimate & Test \\
\midrule
$\hat\beta^{on}_0$ (onset) & $+0.0159$ & margin $+1.59$ pp vs.\ $(0,0)$ \\
$\hat\beta^{per}_0$ (persist) & $-0.0017$ & --- \\
$\hat\beta^{ex}_0$ (exit) & $+0.0109$ & still $+1.09$ pp after clearing \\
Asymmetry $\hat L_0=\hat\beta^{on}_0+\hat\beta^{ex}_0$ & $+0.0268$ & SE $0.0098$, $z=2.73$, $p=0.006$ \\
Dose marginal at $x=5$ (Eq.\ \ref{eq:dose}) & $+0.0023$ & $p=0.010$ \\
Dose marginal at $x=10$ & $+0.0018$ & $p=0.14$ \\
Binned $\hat\Delta=\hat\delta_{10+}-\hat\delta_{1\text{--}3}$ & $+0.018$ & $p=0.033$ \\
\bottomrule
\end{tabular}

\vspace{2pt}
\footnotesize\raggedright
Onset and exit are \emph{both} positive: the margin does not unwind on exit
($H_0$ rejected at 1\%)---asymmetric persistence. The positive dose-response is read
as survivorship/adaptation (financially weaker hospitals exit the unbalanced panel
first), not as climate exposure improving finances; the persistence claim is
anchored to the drought hysteresis result, which needs no long accumulation window.
\end{table}

%================================================================
\section{Paper 3 --- Inequality and provider heterogeneity}
%================================================================
Paper~3 asks whether the effects of Papers~1--2 concentrate in already-disadvantaged
groups. For a moderator $M_{it}$,
\begin{equation}
Y_{it}=\alpha_i+\gamma_t+\beta \Sct+\theta(\Sct\times M_{it})+\phi M_{it}
        +\mathbf X_{it}'\delta+\varepsilon_{it},
\label{eq:het}
\end{equation}
so $\beta$ is the shock effect at $M=0$, $\beta+\theta$ at $M=1$ (SE from
$\sqrt{V_{\beta\beta}+V_{\theta\theta}+2V_{\beta\theta}}$), and $\theta$ the
differential. The demand side uses county social vulnerability (CDC/ATSDR SVI,
top-quartile dummy); the supply side uses hospital safety-net status, ownership,
state Medicaid-expansion, and county hospital-market concentration (HHI). For
strain-increasing outcomes (uncompensated-care \%) $\theta>0$ indicates concentration
in high-$M$ units; for operating margin a negative $\theta$ does.
Table~\ref{tab:het} gives selected interactions.

\begin{table}[h]
\centering
\caption{Heterogeneity: selected shock $\times$ moderator interactions}
\label{tab:het}
\small
\begin{tabular}{llll}
\toprule
Shock $\times$ moderator & Outcome & Direction of $\hat\theta$ & $p$ \\
\midrule
$\textsf{High\_CDD}\times$ SafetyNet & Uncomp.\ care \%NPR & higher for safety-net & $0.021$ \\
Drought $\times$ SVI (county) & Income per capita & larger loss in high-SVI & $<0.05$ \\
$\textsf{High\_HDD}\times$ SafetyNet & Uncomp.\ care \%NPR & higher for non-safety-net & $0.011$ \\
Drought $\times$ Medicaid expansion & Uncomp.\ care \%NPR & higher in expansion states & $<0.001$ \\
\bottomrule
\end{tabular}

\vspace{2pt}
\footnotesize\raggedright
Headline: heat-driven uncompensated care concentrates in safety-net hospitals---the
supply-side analogue of high-SVI county amplification on the demand side. Interactions
are shock- and outcome-specific and not all point the ``vulnerability amplifies'' way;
all are reported transparently.
\end{table}

%================================================================
\section{Inference}
%================================================================
All standard errors cluster at the state level throughout; for HIX-premium outcomes
the county panel additionally reports rating-area--clustered SEs, since counties
sharing a rating area have identical premiums by construction. Every post-estimation
quantity---cumulative lag effects \eqref{eq:cum}, asymmetry statistics \eqref{eq:sym},
dose-level marginal effects \eqref{eq:dose}, and moderator marginal effects
\eqref{eq:het}---is a linear combination $\hat L=\mathbf w'\hat{\boldsymbol\beta}$ with
$\widehat{\operatorname{Var}}(\hat L)=\mathbf w'\hat V\mathbf w$ evaluated on the
state-clustered vcov $\hat V$; tests use the normal approximation and are two-tailed at
$\alpha=0.05$. All models are estimated with \texttt{fixest::feols} (Bergé 2018), which
performs the two-way within transform and exports the clustered vcov for the
lincom steps.

%================================================================
\appendix
\section*{Appendix: estimating-equation reference}
\small
\begin{center}
\begin{tabular}{lll}
\toprule
Paper & Estimating equation & Key quantity \\
\midrule
1 Incidence & $Y_{it}=\alpha_i+\gamma_t+\sum_{k=0}^2\beta_k S_{i,t-k}+\varepsilon$
  & $\hat\beta_{\mathrm{cum}}=\sum_k\hat\beta_k$ \\
1 Incidence (DiD) & $Y_{ct}=\alpha_c+\gamma_t+\tau(\text{Tr}_c{\times}\text{Post}_t)+\varepsilon$
  & ATT $\hat\tau=\Delta^T-\Delta^C$ \\
2 Persistence (onset/exit) & $Y_{i,t+h}=\alpha_i+\gamma_t+\beta^{on}D^{on}+\beta^{per}D^{per}+\beta^{ex}D^{ex}+\varepsilon$
  & $\hat L_h=\hat\beta^{on}_h+\hat\beta^{ex}_h$ \\
2 Persistence (dose) & $Y_{it}=\alpha_i+\gamma_t+\beta_1 C_{it}+\beta_2 C_{it}^2+\varepsilon$
  & $\hat\beta_1+2\hat\beta_2 x$ \\
3 Heterogeneity & $Y_{it}=\alpha_i+\gamma_t+\beta S_{it}+\theta(S_{it}{\times}M_{it})+\phi M_{it}+\varepsilon$
  & differential $\hat\theta$ \\
\bottomrule
\end{tabular}
\end{center}

\vspace{2pt}
\footnotesize\noindent
Estimation in R (\texttt{fixest}). Shock construction:
\texttt{process\_county\_climate.R}, \texttt{process\_hospital\_panel.R}.
Persistence helpers: \texttt{transition\_symmetry.R}, \texttt{cumulative\_dose.R}.
Natural-experiment DiD: \texttt{run\_did\_analysis.R};
doubly-robust check: \texttt{did\_robustness/02\_doubly\_robust\_did.R}.

\end{document}
